\chapterimage{figs/chapterhead.png}
\chapter{Evidence and Maturity Matrix}
\label{chap:appendix-evidence-and-maturity-matrix}

This appendix expands the evidence ladder used in the front matter into a
readable matrix. It is meant to keep the language of the manual consistent:
build success is not release readiness, and startup evidence is not scientific
validation. \index{evidence}\index{maturity}

\begin{table}[htbp]
\centering
\small
\begin{tabular}{|p{2.8cm}|p{5.1cm}|p{4.0cm}|}
\hline
\textbf{Evidence level} & \textbf{What it supports} & \textbf{What it does not support} \\
\hline
Build evidence & The current source tree can be configured and compiled for the selected preset. & It does not prove runtime correctness or scientific validity. \\
\hline
Startup evidence & The executable launches, opens its surface, and shuts down cleanly. & It does not prove complete workflow coverage. \\
\hline
Functional evidence & A focused user flow or API path behaves as documented. & It does not prove all edge cases or all configuration combinations. \\
\hline
Numerical evidence & A bounded computation reproduces the observed output for the documented inputs. & It does not guarantee universal accuracy or stability. \\
\hline
Scientific evidence & The model, reference, and interpretation were reviewed against an explicit domain basis. & It does not follow from a compiler pass alone. \\
\hline
Security evidence & The documented boundary was inspected and a local or private deployment claim was confirmed. & It does not authorize public exposure or production trust. \\
\hline
Release readiness & The documented validation, review, and governance gates were cleared. & It does not arise from one local validation run. \\
\hline
\end{tabular}
\end{table}

The current manual aligns the evidence ladder with the simulation literature
and the repository's validation chapters. \cite{banks2005,law2015,myers2011,sommerville2015}\index{simulation evidence}\index{release readiness}

The practical lesson is simple: each claim must stay inside the evidence level
that actually supports it.
